% =====================================================================
% frontmatter/acronyms.tex
% PT: Lista de siglas e acrónimos. Cada sigla é definida uma vez aqui e
%     usada no texto com \gls{rótulo}. A lista é ordenada automaticamente.
% EN: List of acronyms. Each acronym is defined once here and used in the
%     text with \gls{label}. The list is sorted automatically.
%
% PT: Formato: \newacronym{rótulo}{SIGLA}{Forma por extenso}
%     No texto:
%       \gls{ai}       1.º uso: "Artificial Intelligence (AI)"; depois: "AI"
%       \glspl{cnn}    plural: "Convolutional Neural Networks (CNNs)" / "CNNs"
%       \Gls{ai}       com maiúscula inicial, no início de uma frase
%       \acrshort{ai}  só a sigla: "AI"
%       \acrlong{ai}   só a forma por extenso: "Artificial Intelligence"
%       \acrfull{ai}   sempre completa: "Artificial Intelligence (AI)"
% EN: Format: \newacronym{label}{ACRONYM}{Long form}
%     In the text:
%       \gls{ai}       1st use: "Artificial Intelligence (AI)"; then: "AI"
%       \glspl{cnn}    plural: "Convolutional Neural Networks (CNNs)" / "CNNs"
%       \Gls{ai}       capitalised, at the start of a sentence
%       \acrshort{ai}  acronym only: "AI"
%       \acrlong{ai}   long form only: "Artificial Intelligence"
%       \acrfull{ai}   always full: "Artificial Intelligence (AI)"
% =====================================================================

% PT: Exemplos em inglês.
% EN: English examples.
\newacronym{ai}{AI}{Artificial Intelligence}
\newacronym{ml}{ML}{Machine Learning}
\newacronym{cnn}{CNN}{Convolutional Neural Network}

% PT: Exemplo de sigla portuguesa, para texto escrito em português.
% EN: Example of a Portuguese acronym, for text written in Portuguese.
\newacronym{ia}{IA}{Inteligência Artificial}
